\documentclass[handout,nooutcomes,noauthor]{ximera}

\graphicspath{  
{./}
{./whoAreYou/}
{./drawingWithTheTurtle/}
{./bisectionMethod/}
{./circles/}
{./anglesAndRightTriangles/}
{./lawOfSines/}
{./lawOfCosines/}
{./plotter/}
{./staircases/}
{./pitch/}
{./qualityControl/}
{./symmetry/}
{./nGonBlock/}
}


%% page layout
\usepackage[cm,headings]{fullpage}
\raggedright
\setlength\headheight{13.6pt}


%% fonts
\usepackage{euler}

\usepackage{FiraMono}
\renewcommand\familydefault{\ttdefault} 
\usepackage[defaultmathsizes]{mathastext}
\usepackage[htt]{hyphenat}

\usepackage[T1]{fontenc}
\usepackage[scaled=1]{FiraSans}

%\usepackage{wedn}
\usepackage{pbsi} %% Answer font


\usepackage{cancel} %% strike through in pitch/pitch.tex


%% \usepackage{ulem} %% 
%% \renewcommand{\ULthickness}{2pt}% changes underline thickness

\tikzset{>=stealth}

\usepackage{adjustbox}

\setcounter{titlenumber}{-1}

%% journal style
\makeatletter
\newcommand\journalstyle{%
  \def\activitystyle{activity-chapter}
  \def\maketitle{%
    \addtocounter{titlenumber}{1}%
                {\flushleft\small\sffamily\bfseries\@pretitle\par\vspace{-1.5em}}%
                {\flushleft\LARGE\sffamily\bfseries\thetitlenumber\hspace{1em}\@title \par }%
                {\vskip .6em\noindent\textit\theabstract\setcounter{question}{0}\setcounter{sectiontitlenumber}{0}}%
                    \par\vspace{2em}
                    \phantomsection\addcontentsline{toc}{section}{\thetitlenumber\hspace{1em}\textbf{\@title}}%
                     }}
\makeatother



%% thm like environments
\let\question\relax
\let\endquestion\relax

\newtheoremstyle{QuestionStyle}{\topsep}{\topsep}%%% space between body and thm
		{}                      %%% Thm body font
		{}                              %%% Indent amount (empty = no indent)
		{\bfseries}            %%% Thm head font
		{)}                              %%% Punctuation after thm head
		{ }                           %%% Space after thm head
		{\thmnumber{#2}\thmnote{ \bfseries(#3)}}%%% Thm head spec
\theoremstyle{QuestionStyle}
\newtheorem{question}{}



\let\freeResponse\relax
\let\endfreeResponse\relax

%% \newtheoremstyle{ResponseStyle}{\topsep}{\topsep}%%% space between body and thm
%% 		{\wedn\bfseries}                      %%% Thm body font
%% 		{}                              %%% Indent amount (empty = no indent)
%% 		{\wedn\bfseries}            %%% Thm head font
%% 		{}                              %%% Punctuation after thm head
%% 		{3ex}                           %%% Space after thm head
%% 		{\underline{\underline{\thmname{#1}}}}%%% Thm head spec
%% \theoremstyle{ResponseStyle}

\usepackage[tikz]{mdframed}
\mdfdefinestyle{ResponseStyle}{leftmargin=1cm,linecolor=black,roundcorner=5pt,
, font=\bsifamily,}%font=\wedn\bfseries\upshape,}


\ifhandout
\NewEnviron{freeResponse}{}
\else
%\newtheorem{freeResponse}{Response:}
\newenvironment{freeResponse}{\begin{mdframed}[style=ResponseStyle]}{\end{mdframed}}
\fi



%% attempting to automate outcomes.

%% \newwrite\outcomefile
%%   \immediate\openout\outcomefile=\jobname.oc
%% \renewcommand{\outcome}[1]{\edef\theoutcomes{\theoutcomes #1~}%
%% \immediate\write\outcomefile{\unexpanded{\outcome}{#1}}}

%% \newcommand{\outcomelist}{\begin{itemize}\theoutcomes\end{itemize}}

%% \NewEnviron{listOutcomes}{\small\sffamily
%% After answering the following questions, students should be able to:
%% \begin{itemize}
%% \BODY
%% \end{itemize}
%% }
\usepackage[tikz]{mdframed}
\mdfdefinestyle{OutcomeStyle}{leftmargin=2cm,rightmargin=2cm,linecolor=black,roundcorner=5pt,
, font=\small\sffamily,}%font=\wedn\bfseries\upshape,}
\newenvironment{listOutcomes}{\begin{mdframed}[style=OutcomeStyle]After answering the following questions, students should be able to:\begin{itemize}}{\end{itemize}\end{mdframed}}



%% my commands

\newcommand{\snap}{{\bfseries\itshape\textsf{Snap!}}}
\newcommand{\flavor}{\link[\snap]{https://snap.berkeley.edu/}}
\newcommand{\mooculus}{\textsf{\textbf{MOOC}\textnormal{\textsf{ULUS}}}}


\usepackage{tkz-euclide}
\tikzstyle geometryDiagrams=[rounded corners=.5pt,ultra thick,color=black]
\colorlet{penColor}{black} % Color of a curve in a plot



\ifhandout\newcommand{\mynewpage}{\newpage}\else\newcommand{\mynewpage}{}\fi

\title{Rep-tiles and area}

\author{Bart Snapp}

\begin{document}
\begin{abstract}
  We will investigate more interesting rep-tiles.
\end{abstract}
\maketitle


\begin{listOutcomes}
\item Demonstrate given shapes are rep-n-tiles.
\item Understand a relationship between scaling, perimeter, and area.
\end{listOutcomes}



Figureing out \textbf{how} shapes are rep-tiles can be difficult if
you don't actually have phyical copies of the rep-tiles to play with.
I suggest you have someone help you write a \snap\ script using the block:
\begin{center}
  \includegraphics{BLOCKreptile.png}
\end{center}
I'll explain the arguments (blanks in the BLOCK) for completness sake.
Working right-to-left, the first TURTLE BLANK is for a costume, the
second TURTLE BLANK is for a horizontal reflection of the costume, the
blank labled $r:$ is the size of the rotation in degrees, and the
blank labeled $w:$ is the desired width of the rep-tile.
\begin{itemize}
\item The ARROW keys will move the rep-tile around.
\item Pressing $r$ will rotate the rep-tile clockwise the specified
  number of degrees.
\item Pressing $f$ will flip the turtle horizontally across a line
  perpendicular to its base.
\end{itemize}


\mynewpage

\begin{question}
Consider the following polygons:
\[
\includegraphics[width=.2\textwidth]{rhomb.png} \qquad
\includegraphics[width=.15\textwidth]{P.png}\qquad
\includegraphics[width=.2\textwidth]{wedge.png} 
\]
Use the REPTILE BLOCK to show that each of these polygons is a
rep-$4$-tile. A SCREENSHOT of the \snap\ STAGE is sufficent.
\end{question}

\mynewpage


\begin{question}
  Now we'll investigate a rep-$4$-tile and a rep-$5$-tile.
  \begin{enumerate}
  \item Show that the following polygon is a rep-$4$-tile:
    \begin{center}
      \includegraphics[width=.2\textwidth]{sphinx.png}
    \end{center}
  \item Show that a $(1,2,\sqrt{5})$ right triangle is a
    rep-$5$-tile.
    \begin{hint}
      I suggest you rotate $90^\circ$.
    \end{hint}
  \end{enumerate}
  In each case, if you want to convince me that you have a rep-tile,
  you must EXPLAIN why the angles are unchanged, and why the
  proportions of the linear measurements are unchanged.
\end{question}

\mynewpage

\begin{question}
  Now we'll connect rep-tiles to SCALING area.
  \begin{enumerate}
    \item EXPLAIN why given ANY rep-2-tile,
      \[
      \frac{\text{new perimeter}}{\text{old perimeter}}\qquad\text{and}\qquad \frac{\text{new area}}{\text{old area}}
      \]
      will be constants that are invariant, regardless of the shape of
      the given rep-2-tile.
    \item EXPLAIN why given ANY rep-$n$-tile,
      \[
      \frac{\text{new perimeter}}{\text{old perimeter}}\qquad\text{and}\qquad \frac{\text{new area}}{\text{old area}}
      \]
      will be constants that are invariant, regardless of the shape of
      the given rep-$n$-tile.
    \item
      Start make a table like the one below:
      \begin{center}
        \begin{tabular}{c|c|c}
          rep-$n$-tile & $\frac{\text{new perimeter}}{\text{old perimeter}}$ & $\frac{\text{new area}}{\text{old area}}$  \\
          \hline\hline
          2 & $\vdots$  &  $\vdots$  \\ 
          3 & $\vdots$  &  $\vdots$  \\ 
        \end{tabular}
      \end{center}
      COMPLETE the table up to $n= 10$.
  \end{enumerate}
\end{question}



\end{document}
